\section{Introduction}

Reversing an articulated tractor--semitrailer is a nonlinear and nonholonomic
control problem. In reverse motion, small steering errors can increase the
relative heading between the tractor and trailer, causing tracking failure,
collision, or jackknifing. Loading-dock operation adds tight lateral clearance
and final-pose requirements, making it a useful setting for comparing learned
and model-based control.

The project studies whether a PPO policy using trajectory-relative tracking
errors and preview information can reliably execute reverse docking and how its
behavior compares with NMPC when initial pose, trailer geometry, parking-bay
dimensions, and state or model uncertainty change. The comparison is not
designed to make one controller win. It is intended to characterize accuracy,
safety, computation, robustness, and transfer behavior under a common task
definition.

The required contribution of the course project is an implemented and
quantitatively evaluated two-dimensional system. A three-dimensional MuJoCo
model is a secondary integration target. If full transfer is not reliable or
cannot be completed within the schedule, the complete controller comparison
will remain in the verified two-dimensional environment, while MuJoCo will be
used for a smaller validated demonstration set.

This is a living report. Statements are limited to implemented behavior and
saved evidence. Planned methods and experiments are labeled as such, and no
novelty or performance claim is made before comparative results exist.
